% Generated draft: requires mathematical and source review.
\documentclass[11pt]{article}
\usepackage[margin=25mm]{geometry}
\usepackage{fontspec}
\setmainfont{Latin Modern Roman}
\usepackage{amsmath,amssymb,mathtools,unicode-math}
\setmathfont{Latin Modern Math}
\usepackage{enumitem,needspace,xurl,hyperref}
\setlength{\parindent}{0pt}
\setlength{\parskip}{5pt}
\setlength{\emergencystretch}{3em}
\newcommand{\N}{\mathbb{N}}
\newcommand{\Z}{\mathbb{Z}}
\newcommand{\Q}{\mathbb{Q}}
\newcommand{\R}{\mathbb{R}}
\newcommand{\C}{\mathbb{C}}
\newcommand{\F}{\mathbb{F}}
\newcommand{\A}{\mathbb{A}}
\newcommand{\PP}{\mathbb{P}}
\newcommand{\E}{\mathbb{E}}
\newcommand{\eps}{\varepsilon}
\newcommand{\dd}{\,\mathrm{d}}
\begin{document}
\section*{UnsolvedMath ID: 701363}
\textbf{Source problem code:} catalogue entry
\subsection{Definitions and mathematical statement}
A group \(G\) is finitely generated if there is a finite subset \(X\subseteq G\) that generates \(G\). It is finitely presented if there exist a finite set \(X\) and a finite subset \(R\) of the free group \(F(X)\) such that
\[
G\cong F(X)/\langle\!\langle R\rangle\!\rangle,
\]
where \(\langle\!\langle R\rangle\!\rangle\) denotes the normal subgroup of \(F(X)\) generated by \(R\).

The supplied record uses the notation \(\mathrm{CT}_{\mathcal P}(\Z)\) but does not define this group or specify the admissible sets \(\mathcal P\). Consequently, the intended question can only be stated relative to these unspecified source definitions:
\begin{quote}
For every admissible finite set \(\mathcal P\), is the group \(\mathrm{CT}_{\mathcal P}(\Z)\) finitely presented?
\end{quote}
Equivalently, for each such \(\mathcal P\), do there exist a finite set \(X\) and a finite subset \(R\subseteq F(X)\) for which
\[
\mathrm{CT}_{\mathcal P}(\Z)\cong F(X)/\langle\!\langle R\rangle\!\rangle?
\]
No uniform choice or bound for the generators or relations is requested.

As background, the supplied record asserts, for the intended admissible parameters, that \(\mathrm{CT}_{\mathcal P}(\Z)\) is finitely generated if and only if \(\mathcal P\) is finite. It further reports that \(\mathrm{CT}_{\varnothing}(\Z)\) is isomorphic to the finitely presented first Higman--Thompson group, attributing this observation to J. P. McDermott and referring to Remark 1.4 of the cited 2017 paper by S. Kohl. These assertions are retained as supplied background. The missing definition of \(\mathrm{CT}_{\mathcal P}(\Z)\) and the missing domain of \(\mathcal P\) prevent a self-contained formulation.
\subsection{Short English statement}
S. Kohl asks whether \(\mathrm{CT}_{\mathcal P}(\Z)\) is finitely presented for every finite admissible set \(\mathcal P\). The supplied background reports finite generation exactly when \(\mathcal P\) is finite and finite presentability for \(\mathcal P=\varnothing\). The group construction and the allowed parameters are not defined in the record.
\subsection{Sources}
\begin{itemize}
\item Supplied local UnsolvedMath registry record: permanent ID \(701363\), exact source problem code \texttt{catalogue entry}, title \emph{Kourovka Notebook Problem 17.61}. The registry labels the problem \texttt{open}; this is supplied metadata, not an independently verified status assessment.
\item \emph{The Kourovka Notebook (Unsolved Problems in Group Theory)}, original problem 17.61, proposed by S. Kohl in the 17th issue (2010). The supplied background also identifies the 21st issue edition (2026), \texttt{arXiv:1401.0300v48}. The supplied collection description names E. I. Khukhro and V. D. Mazurov as editors. These source identifications have not been independently checked.
\item S. Kohl, \emph{J. Group Theory} \textbf{20}, no. 5 (2017), 1025--1030, Remark 1.4, cited in the supplied statement for the observation attributed to J. P. McDermott concerning the empty-parameter case. The record supplies neither the article title nor a URL.
\end{itemize}
\end{document}
